\vspace{-0.2cm}
\section{Preliminaries and Problem Setup}
\label{sec:prelim}
\vspace{-0.2cm}

Our goal is to develop an approach for training LLMs to improve their own predictions entirely on self-generated data. As discussed so far, we situate ourselves in the intrinsic self-correction setting~\citep{huang2023large}, where models attempt to correct their initial responses \emph{without} any external feedback. Concretely, given a dataset $\mathcal{D} = \{(\bx_i, \by^*_i) \}_{i=1}^N$ of problems $\bx_i$ and responses $\by^*_i$, we will train an LLM policy $\pi_\theta(\cdot|[\bx, \hat{\by}_{1:l}, p_{1:l}])$ that, given the problem $\bx$, previous $l$ model attempts $\hat{\by}_{1:l}$ at the problem, and auxiliary instructions $p_{1:l}$ (e.g., instruction to find a mistake and improve the response), solves the problem $\bx$ as correctly as possible. This formalism is akin to the multi-turn MDP in \citet{qu2024recursive}. We also assume access to an oracle reward $\widehat{r}(\by, \by^*)$, such as an answer checker~\citep{uesato2022solving}, that evaluates the correctness of response $\by$ by comparing it with the oracle response $\by^*$. Critically, we \emph{do not} assume access to this oracle at test-time; instead, the model must deduce whether there was a mistake and correct it if necessary, as is often the case in e.g. mathematical reasoning problems. Unlike the setup of \citet{qu2024recursive}, we also do not run majority voting for most of our main results. Two example traces of self-correction are given in Figure~\ref{fig:score_main}, and our problem setting is depicted pictorially in Figure~\ref{fig:problem_setup}.

We aim to find an LLM policy $\pi(\square|\circ)$ mapping input tokens $\circ$ to output tokens $\square$ that maximizes the correctness reward obtained from the verifier at the end of $l+1$ turns ($l = 1$). Formally:
\begin{align}
\label{eq:prelims_correctness}
    \max_{\pi_\theta}~~ \mathbb{E}_{\bx, \by^* \sim \mathcal{D}, \hat{\by}_{l+1} \sim \pi_\theta \left(\cdot|[\bx, \hat{\by}_{1:l}, p_{1:l}]\right)} \left[ \sum_{i=1}^{l+1} \widehat{r} \left( \hat{\by}_{i}, \by^* \right) \right].
\end{align}
Crucially, note that unlike standard SFT or prevalent RL fine-tuning workflows, which train the policy $\pi$ to directly produce $\by^*$ (or any other $\by$ wih $\widehat{r}(\by, \by^*) = 1$), Equation~\ref{eq:prelims_correctness} trains $\pi$ over multiple attempts \emph{simultaneously}, where intermediate turns are supervised indirectly to maximize the sum. 

\textbf{Base RL fine-tuning approach we use.} We use a REINFORCE policy gradient training approach with a KL-divergence penalty against a fixed model~\citep{ahmadian2024back}, which is widely used in RL fine-tuning of LLMs, primarily in the setting of single-turn RLHF. Formally, these methods train the policy $\pi_{\theta}(\cdot|\bx)$ to optimize the following, where $\pi_\mathrm{ref}$ is a reference policy.
\begin{align}
    \label{eq:base_reinforce}
    \max_{\theta}~~ \expected_{\bx_t, \by_t \sim \pi_{\theta}(\cdot | \bx_t)}\left[\widehat{r}(\by_t, \by^*) - \beta_1 D_{KL}(\pi_\theta(\cdot | \bx_t) || \pi_\mathrm{ref}(\cdot | \bx_t))\right],
\end{align}
\textbf{Metrics.} To measure self-correction performance (we consider $l=2$ in this paper), we report and analyze the following metrics: \textbf{(1)} \textbf{Accuracy@t1}: the model's accuracy at the first attempt; \textbf{(2)} \textbf{Accuracy@t2}: the model's accuracy at the second attempt, \textbf{(3)} \textbf{$\Delta$(t1, t2)}: the net improvement in model accuracy between the first and second attempts, which measures the efficacy of self-correction, \textbf{(4)} \textbf{$\Delta^{\mathrm{i} \rightarrow \mathrm{c}}$(t1, t2)}: the fraction of problems that are incorrect in the first attempt but become correct at the second attempt, which measures how many \emph{new} problems can self-correction solve; and \textbf{(5)} \textbf{$\Delta^{\mathrm{c} \rightarrow \mathrm{i}}$(t1, t2)}: the fraction of problems that are correct in the first attempt but become incorrect at the second attempt, which measures how well the model understands what makes a response correct.




